% Erratum E13 for inclusion in main.tex.
% Suggested use: % Erratum E13 for inclusion in main.tex.
% Suggested use: % Erratum E13 for inclusion in main.tex.
% Suggested use: % Erratum E13 for inclusion in main.tex.
% Suggested use: \input{erratum_E13_landweber_qok_nu0}

\subsection*{Erratum E13: the parameter range in equation (67)}
\label{s:erratum:E13}

The statement surrounding equation (67) should require $\nu>0$, or else
an explicit limiting convention should be imposed.  As printed, the factor
\[
        (1+k/\nu)^\nu
\]
is undefined when $\nu=0$.  Thus the sentence following equation (67)
should read that the hidden constants may depend on the fixed parameter
$\nu>0$.

This is only a parameter-range correction.  It is consistent with
Proposition~3.4, which is also stated for $\nu>0$.  Indeed, when
$r=0$ the normalization in Proposition~3.4 is
\[
        \varepsilon_k
        =
        \frac{1}{\sigma}\frac{\nu}{k+\nu},
\]
and therefore
\[
        \varepsilon_k^{-\nu}
        =
        \sigma^\nu(1+k/\nu)^\nu .
\]
Since $\sigma$ is fixed, the factor in equation (67) is just the
Landweber rate normalization from Proposition~3.4, up to the
constant $\sigma^\nu$.  This normalization is meaningful for $\nu>0$.

If one wants to include the endpoint $\nu=0$ formally, the natural
convention is the limit
\[
        \lim_{\nu\downarrow0}(1+k/\nu)^\nu = 1
        \qquad (k \text{ fixed}),
\]
and then equation (67) becomes the boundedness comparison
\[
        \sup_{\norm{ y^\dag - y^\delta }{} \le \delta}
        \norm{x^\dag - x_{\bar{k}(\delta,y^\delta,\ldots)}^\delta}{}
        \sim
        \sup_{k\geq0}
        \norm{x^\dag - x_k}{} .
\]
This endpoint convention, however, is separate from the algebraic-rate
case treated in Proposition~3.4.  The simplest correction is
therefore to state equation (67) only for $\nu>0$.


\subsection*{Erratum E13: the parameter range in equation (67)}
\label{s:erratum:E13}

The statement surrounding equation (67) should require $\nu>0$, or else
an explicit limiting convention should be imposed.  As printed, the factor
\[
        (1+k/\nu)^\nu
\]
is undefined when $\nu=0$.  Thus the sentence following equation (67)
should read that the hidden constants may depend on the fixed parameter
$\nu>0$.

This is only a parameter-range correction.  It is consistent with
Proposition~3.4, which is also stated for $\nu>0$.  Indeed, when
$r=0$ the normalization in Proposition~3.4 is
\[
        \varepsilon_k
        =
        \frac{1}{\sigma}\frac{\nu}{k+\nu},
\]
and therefore
\[
        \varepsilon_k^{-\nu}
        =
        \sigma^\nu(1+k/\nu)^\nu .
\]
Since $\sigma$ is fixed, the factor in equation (67) is just the
Landweber rate normalization from Proposition~3.4, up to the
constant $\sigma^\nu$.  This normalization is meaningful for $\nu>0$.

If one wants to include the endpoint $\nu=0$ formally, the natural
convention is the limit
\[
        \lim_{\nu\downarrow0}(1+k/\nu)^\nu = 1
        \qquad (k \text{ fixed}),
\]
and then equation (67) becomes the boundedness comparison
\[
        \sup_{\norm{ y^\dag - y^\delta }{} \le \delta}
        \norm{x^\dag - x_{\bar{k}(\delta,y^\delta,\ldots)}^\delta}{}
        \sim
        \sup_{k\geq0}
        \norm{x^\dag - x_k}{} .
\]
This endpoint convention, however, is separate from the algebraic-rate
case treated in Proposition~3.4.  The simplest correction is
therefore to state equation (67) only for $\nu>0$.


\subsection*{Erratum E13: the parameter range in equation (67)}
\label{s:erratum:E13}

The statement surrounding equation (67) should require $\nu>0$, or else
an explicit limiting convention should be imposed.  As printed, the factor
\[
        (1+k/\nu)^\nu
\]
is undefined when $\nu=0$.  Thus the sentence following equation (67)
should read that the hidden constants may depend on the fixed parameter
$\nu>0$.

This is only a parameter-range correction.  It is consistent with
Proposition~3.4, which is also stated for $\nu>0$.  Indeed, when
$r=0$ the normalization in Proposition~3.4 is
\[
        \varepsilon_k
        =
        \frac{1}{\sigma}\frac{\nu}{k+\nu},
\]
and therefore
\[
        \varepsilon_k^{-\nu}
        =
        \sigma^\nu(1+k/\nu)^\nu .
\]
Since $\sigma$ is fixed, the factor in equation (67) is just the
Landweber rate normalization from Proposition~3.4, up to the
constant $\sigma^\nu$.  This normalization is meaningful for $\nu>0$.

If one wants to include the endpoint $\nu=0$ formally, the natural
convention is the limit
\[
        \lim_{\nu\downarrow0}(1+k/\nu)^\nu = 1
        \qquad (k \text{ fixed}),
\]
and then equation (67) becomes the boundedness comparison
\[
        \sup_{\norm{ y^\dag - y^\delta }{} \le \delta}
        \norm{x^\dag - x_{\bar{k}(\delta,y^\delta,\ldots)}^\delta}{}
        \sim
        \sup_{k\geq0}
        \norm{x^\dag - x_k}{} .
\]
This endpoint convention, however, is separate from the algebraic-rate
case treated in Proposition~3.4.  The simplest correction is
therefore to state equation (67) only for $\nu>0$.


\subsection*{Erratum E13: the parameter range in equation (67)}
\label{s:erratum:E13}

The statement surrounding equation (67) should require $\nu>0$, or else
an explicit limiting convention should be imposed.  As printed, the factor
\[
        (1+k/\nu)^\nu
\]
is undefined when $\nu=0$.  Thus the sentence following equation (67)
should read that the hidden constants may depend on the fixed parameter
$\nu>0$.

This is only a parameter-range correction.  It is consistent with
Proposition~3.4, which is also stated for $\nu>0$.  Indeed, when
$r=0$ the normalization in Proposition~3.4 is
\[
        \varepsilon_k
        =
        \frac{1}{\sigma}\frac{\nu}{k+\nu},
\]
and therefore
\[
        \varepsilon_k^{-\nu}
        =
        \sigma^\nu(1+k/\nu)^\nu .
\]
Since $\sigma$ is fixed, the factor in equation (67) is just the
Landweber rate normalization from Proposition~3.4, up to the
constant $\sigma^\nu$.  This normalization is meaningful for $\nu>0$.

If one wants to include the endpoint $\nu=0$ formally, the natural
convention is the limit
\[
        \lim_{\nu\downarrow0}(1+k/\nu)^\nu = 1
        \qquad (k \text{ fixed}),
\]
and then equation (67) becomes the boundedness comparison
\[
        \sup_{\norm{ y^\dag - y^\delta }{} \le \delta}
        \norm{x^\dag - x_{\bar{k}(\delta,y^\delta,\ldots)}^\delta}{}
        \sim
        \sup_{k\geq0}
        \norm{x^\dag - x_k}{} .
\]
This endpoint convention, however, is separate from the algebraic-rate
case treated in Proposition~3.4.  The simplest correction is
therefore to state equation (67) only for $\nu>0$.
